%region Introducción
\newpage
\pagenumbering{arabic}
\setcounter{page}{1}
\section{Introducción}
El reconocimiento automático de matrículas vehiculares, conocido por sus siglas en inglés como LPR (\textit{License Plate Recognition}) o ALPR (\textit{Automatic License Plate Recognition}), ha tomado relevancia dentro del ecosistema de sistemas inteligentes de transporte y videovigilancia \cite{chang-chen2004}. Su objetivo principal consiste en extraer la cadena alfanumérica que identifica a un vehículo a partir de una imagen capturada por una cámara convencional (color, blanco y negro o infrarroja), combinando para ello técnicas de detección de objetos, procesamiento digital de imágenes y reconocimiento óptico de caracteres (OCR) \cite{bohou-li2022}. Un sistema LPR robusto puede aplicarse en control de acceso a estacionamientos \cite{sirithinaphong-chamnongthai1999}, peajes automáticos \cite{davies-emmott1990}, verificación vehicular por cuerpos de seguridad, gestión de tráfico y cobro electrónico de infracciones \cite{yamaguchi-nagaya1999}, entre otros.

En los últimos años, impulsado por la creciente popularidad de las redes neuronales profundas y los modelos basados en transformadores de visión (\textit{Vision Transformer}), el entrenamiento e implementación de sistemas LPR se ha generalizado \cite{zherzdev-gruzdev2018, jain-sasindran2016, khan-ilyas2023}. No obstante, su desarrollo aún enfrenta desafíos. La variabilidad de las condiciones de captura (iluminación variable, ángulos extremos, condiciones climáticas adversas), la diversidad de diseños y formatos de matrículas entre distintas jurisdicciones, y las restricciones legales y éticas asociadas al uso de imágenes reales de vehículos y sus propietarios constituyen barreras que dificultan la obtención de conjuntos de datos suficientemente diversos y etiquetados \cite{du-ibrahim2013}.

En el caso particular de México, esta problemática se ve agravada por la existencia de 32 entidades federativas, cada una con libertad para emplear combinaciones de colores, tipografías y elementos de seguridad personalizados, siempre que se respete la Norma Oficial Mexicana NOM-001-SCT-2-2016 \cite{NOM001SCT22016} en cuanto a dimensiones, zonas de seguridad y caracteres obligatorios. Esta heterogeneidad, sumada a la ausencia de bases de datos públicas etiquetadas que reflejen dicha diversidad, limita severamente la capacidad de investigación y desarrollo de nuevos sistemas de reconocimiento de matrículas mexicanas.

Ante este vacío, se introduce LPR-Sentinel, un sistema completo de generación sintética, aumentación, detección y reconocimiento de matrículas vehiculares mexicanas, enfocado exclusivamente en vehículos particulares. El sistema propuesto se integra por cinco módulos que operan de manera secuencial: MLP-Generator (generación procedimental de matrículas sintéticas conforme a la NOM-001-SCT-2-2016), MLP-Augmentator (aplicación de transformaciones geométricas, ópticas y morfológicas), MLP-Detector (localización de la placa sobre imágenes completas), MLP-Recognizer (reconocimiento de caracteres a través de un modelo OCR) y MLP-Pipeline (integración de todos los componentes y consultas a una base de datos). Una de las contribuciones más relevantes del trabajo es la adopción del formato ONNX (\textit{Open Neural Network Exchange}) como interfaz unificada de despliegue, permitiendo desacoplar el entrenamiento (realizado en entornos de escritorio con TensorFlow y PyTorch) de la inferencia (ejecutada en dispositivos de bajo consumo energético y computacional). Todo el código fuente del sistema se encontrará disponible en \cite{murillolog2026} tras la validación de los resultados.

El resto del trabajo se estructura de la siguiente manera: en la Sección II se describe en extenso cada uno de los cinco módulos que componen al sistema LPR-Sentinel. La Sección III expone los experimentos y resultados obtenidos en conjunto con la descripción del entorno de desarrollo y los requisitos de hardware y software para su replicación. La Sección IV discute los hallazgos más relevantes del trabajo y presenta las conclusiones y líneas de trabajo futuro.
%endregion